%! TEX root = hukum_temu1.tex
% the main TeX file which is intended to compile, :VimtexReload after adjustment
   % this file should be included in the main file
% :h vimtex-tex-root
\documentclass[../main.tex]{subfiles}
% ensure the class of docs
% enumerate may not work


\title{DASAR HUKUM K3 KONSTRUKSI DAN SISTEM MANAJEMEN K3}
\author{Nama Pemateri}
\date{\today}

\begin{document}

% ---------------------------------------------------
% FRAME 1: Judul
% ---------------------------------------------------
\begin{frame}[mytitle=center,c,mycolor=digiPH_ocean!15]
	\frametitle{DASAR HUKUM K3 KONSTRUKSI}
	\framesubtitle{Dan Sistem Manajemen K3}

	\begin{block}{Fokus Pembahasan:}
		\begin{itemize}
			\item Kerangka hukum K3 di Indonesia
			\item K3 pada pembangunan jalan dan jembatan
			\item Kewajiban pengurus dan penyedia jasa
			\item Sistem Manajemen Keselamatan dan Kesehatan Kerja (SMK3)
			\item Siklus penerapan SMK3
		\end{itemize}
	\end{block}

	\vspace{0.5cm}
	\textbf{Kata kunci:}\\
	K3 \textbullet\ Konstruksi \textbullet\ Jalan \textbullet\ Jembatan \textbullet\ SMK3 \textbullet\ SMKK \textbullet\ Manajemen Risiko
\end{frame}

% ---------------------------------------------------
% FRAME 2: Capaian Pembelajaran
% ---------------------------------------------------
\begin{frame}[mytitle=standard,t,mycolor=digiPH_cream!40]
	\frametitle{Capaian Pembelajaran}

	Setelah mengikuti pembelajaran, mahasiswa mampu:
	\begin{itemize}
		\item Mengidentifikasi regulasi K3 yang relevan pada pekerjaan konstruksi jalan dan jembatan.
		\item Menjelaskan ruang lingkup K3 dan kewajiban pengurus tempat kerja.
		\item Menganalisis risiko K3 pada tahapan pekerjaan konstruksi.
		\item Menjelaskan konsep, tujuan, manfaat, dan prinsip SMK3.
		\item Menguraikan lima tahap penerapan SMK3 secara sistematis.
		\item Menghubungkan persyaratan hukum dengan praktik pengendalian K3 di proyek.
	\end{itemize}

	\begin{exampleblock}{Orientasi Pembelajaran:}
		Dari \textbf{“apa dasar hukumnya?”} $\rightarrow$ \textbf{“apa kewajibannya?”} $\rightarrow$ \textbf{“bagaimana diterapkan?”}
	\end{exampleblock}
\end{frame}

% ---------------------------------------------------
% FRAME 3: Mengapa K3 Menjadi Persyaratan Fundamental?
% ---------------------------------------------------
\begin{frame}[mytitle=standard,c,mycolor=white]
	\frametitle{Mengapa K3 Menjadi Persyaratan Fundamental?}

	\begin{block}{Pekerjaan konstruksi memiliki karakteristik:}
		\begin{itemize}
			\item Lingkungan kerja dinamis dan berubah-ubah.
			\item Melibatkan manusia, alat berat, material, energi, dan pekerjaan berisiko tinggi.
			\item Berinteraksi dengan masyarakat dan lalu lintas pada proyek jalan.
			\item Memiliki potensi kecelakaan, penyakit akibat kerja, kerusakan aset, dan gangguan lingkungan.
		\end{itemize}
	\end{block}

	Kesalahan pengendalian dapat berdampak pada tenaga kerja, pengguna jalan, masyarakat, biaya, mutu, waktu, dan keberlangsungan proyek.

	\begin{alertblock}{Prinsip Utama}
		Keselamatan bukan sekadar perlengkapan APD, tetapi sistem pengendalian risiko yang dirancang sejak tahap perencanaan.
	\end{alertblock}
\end{frame}

% ---------------------------------------------------
% FRAME 4a: Hierarki Dasar Hukum K3 Konstruksi (Bagian 1)
% ---------------------------------------------------
\begin{frame}[mytitle=center,t,mycolor=digiPH_gray!30]
	\frametitle{Hierarki Dasar Hukum K3 Konstruksi (1/2)}

	\begin{block}{A. Tingkat dasar}
		\textbf{UUD NRI Tahun 1945}\\
		$\rightarrow$ Menjadi landasan konstitusional perlindungan dan kesejahteraan manusia.
	\end{block}

	\begin{block}{B. Regulasi K3}
		\textbf{UU No. 1 Tahun 1970 tentang Keselamatan Kerja}\\
		$\rightarrow$ Fondasi umum keselamatan kerja dan perlindungan tenaga kerja. Statusnya masih berlaku.
	\end{block}

	\begin{block}{C. Regulasi jasa konstruksi}
		\textbf{UU No. 2 Tahun 2017 tentang Jasa Konstruksi}\\
		$\rightarrow$ Mengatur penyelenggaraan jasa konstruksi termasuk keamanan, keselamatan, kesehatan, dan keberlanjutan konstruksi.
	\end{block}
\end{frame}

% ---------------------------------------------------
% FRAME 4b: Hierarki Dasar Hukum K3 Konstruksi (Bagian 2)
% ---------------------------------------------------
\begin{frame}[mytitle=center,t,mycolor=digiPH_gray!30]
	\frametitle{Hierarki Dasar Hukum K3 Konstruksi (2/2)}

	\begin{block}{D. Regulasi sistem manajemen}
		\textbf{PP No. 50 Tahun 2012}\\
		$\rightarrow$ Mengatur Penerapan Sistem Manajemen Keselamatan dan Kesehatan Kerja. Status berlaku.
	\end{block}

	\begin{block}{E. Regulasi sektoral konstruksi}
		\textbf{Permen PUPR No. 10 Tahun 2021}\\
		$\rightarrow$ Pedoman Sistem Manajemen Keselamatan Konstruksi (SMKK).
	\end{block}
\end{frame}

% ---------------------------------------------------
% FRAME 5: UU No. 1 Tahun 1970 tentang Keselamatan Kerja
% ---------------------------------------------------
\begin{frame}[mytitle=standard,c,mycolor=digiPH_creamy!20]
	\frametitle{UU No. 1 Tahun 1970 tentang Keselamatan Kerja}

	\begin{block}{Pokok Pemikiran UU No. 1 Tahun 1970}
		UU ini merupakan fondasi utama hukum keselamatan kerja di Indonesia. Dalam konteks konstruksi, substansinya menegaskan pentingnya:
		\begin{itemize}
			\item perlindungan tenaga kerja;
			\item pencegahan dan pengendalian kecelakaan;
			\item pengamanan tempat kerja;
			\item pengendalian mesin, peralatan, bahan, dan proses kerja;
			\item penyediaan kondisi kerja yang aman.
		\end{itemize}
	\end{block}

	\begin{alertblock}{Implikasi bagi konstruksi}
		Pengurus tidak cukup hanya menanggapi kecelakaan, tetapi wajib menciptakan kondisi kerja yang memungkinkan kecelakaan dapat dicegah dan risikonya dikendalikan.
	\end{alertblock}
\end{frame}

% ---------------------------------------------------
% FRAME 6: Ruang Lingkup K3 pada Pekerjaan Konstruksi
% ---------------------------------------------------
\begin{frame}[mytitle=center,t,mycolor=white]
	\frametitle{Ruang Lingkup K3 pada Pekerjaan Konstruksi}

	K3 tidak terbatas pada pekerjaan di dalam bangunan. Ruang kerja dapat meliputi:
	\begin{itemize}
		\item \textbf{DARAT}: Pekerjaan tanah, perkerasan, drainase, basecamp.
		\item \textbf{TANAH}: Galian, timbunan, pondasi, utilitas bawah tanah.
		\item \textbf{AIR}: Pekerjaan jembatan, sungai, pekerjaan bawah air.
		\item \textbf{UDARA}: Pekerjaan pada struktur tinggi, pengangkatan material, pekerjaan dengan alat angkat.
		\item \textbf{KETINGGIAN}: Scaffolding, struktur jembatan, erection, pemasangan komponen.
	\end{itemize}

	\vspace{0.3cm}
	\begin{exampleblock}{Prinsip}
		Di mana terdapat aktivitas kerja dan potensi bahaya, di situ terdapat kewajiban pengendalian K3.
	\end{exampleblock}
\end{frame}

% ---------------------------------------------------
% FRAME 7: UU No. 2 Tahun 2017: K3 dalam Jasa Konstruksi
% ---------------------------------------------------
\begin{frame}[mytitle=standard,t,mycolor=digiPH_ocean!10]
	\frametitle{UU No. 2 Tahun 2017: K3 dalam Jasa Konstruksi}

	UU No. 2 Tahun 2017 secara eksplisit memasukkan keamanan, keselamatan, kesehatan, dan keberlanjutan konstruksi sebagai bagian penting penyelenggaraan jasa konstruksi.

	\begin{block}{Pasal 59 — Prinsip Penting}
		Dalam penyelenggaraan jasa konstruksi, Pengguna Jasa dan Penyedia Jasa wajib memenuhi Standar Keamanan, Keselamatan, Kesehatan, dan Keberlanjutan.
	\end{block}

	Standar tersebut sekurang-kurangnya mencakup:
	\begin{itemize}
		\item mutu bahan \& mutu peralatan;
		\item keselamatan dan kesehatan kerja;
		\item prosedur pelaksanaan \& mutu hasil pekerjaan;
		\item operasi dan pemeliharaan;
		\item perlindungan sosial tenaga kerja \& pengelolaan lingkungan.
	\end{itemize}
\end{frame}

% ---------------------------------------------------
% FRAME 8: Dasar Hukum Pembangunan Jalan dan Jembatan
% ---------------------------------------------------
\begin{frame}[mytitle=standard,c,mycolor=digiPH_yellow!15]
	\frametitle{Dasar Hukum Pembangunan Jalan dan Jembatan}

	\begin{block}{Regulasi Utama Bidang Jalan}
		\begin{itemize}
			\item \textbf{UU No. 38 Tahun 2004 tentang Jalan}\\
			      $\rightarrow$ menjadi dasar pengaturan penyelenggaraan jalan.
			\item \textbf{UU No. 2 Tahun 2022}\\
			      $\rightarrow$ perubahan kedua atas UU No. 38 Tahun 2004 tentang Jalan.
			\item \textbf{PP No. 34 Tahun 2006 tentang Jalan}\\
			      $\rightarrow$ mengatur lebih lanjut penyelenggaraan jalan.
		\end{itemize}
	\end{block}

	\begin{alertblock}{Dalam Konteks K3}
		Regulasi jalan harus dibaca bersama dengan: UU K3, UU Jasa Konstruksi, PP SMK3, Pedoman SMKK, serta standar teknis dan keselamatan lalu lintas yang relevan.
	\end{alertblock}
\end{frame}

% ---------------------------------------------------
% FRAME 9: Risiko K3 pada Pembangunan Jalan dan Jembatan
% ---------------------------------------------------
\begin{frame}[mytitle=center,t,mycolor=white]
	\frametitle{Risiko K3 pada Pembangunan Jalan dan Jembatan}

	\begin{columns}[t]
		\begin{column}{0.5\textwidth}
			\textbf{Contoh bahaya pada proyek jalan:}
			\begin{itemize}
				\item tertabrak atau terserempet kendaraan;
				\item alat berat terguling;
				\item pekerja tertimbun saat galian;
				\item paparan debu dan kebisingan;
				\item material panas/asphalt;
				\item bahaya listrik \& kerja malam hari.
			\end{itemize}
		\end{column}

		\begin{column}{0.5\textwidth}
			\textbf{Contoh pada proyek jembatan:}
			\begin{itemize}
				\item jatuh dari ketinggian;
				\item benda jatuh;
				\item lifting dan erection;
				\item tenggelam;
				\item runtuhnya struktur sementara;
				\item pekerjaan di atas lalu lintas/air.
			\end{itemize}
		\end{column}
	\end{columns}

	\vspace{0.5cm}
	\begin{exampleblock}{Pendekatan}
		Bahaya $\rightarrow$ Risiko $\rightarrow$ Pengendalian $\rightarrow$ Verifikasi efektivitas
	\end{exampleblock}
\end{frame}

% ---------------------------------------------------
% FRAME 10: Kewajiban Pengurus Tempat Kerja
% ---------------------------------------------------
\begin{frame}[mytitle=standard,t,mycolor=digiPH_cream!30]
	\frametitle{Kewajiban Pengurus Tempat Kerja}

	\textbf{Pengurus harus memastikan:}
	\begin{itemize}
		\item tersedia kebijakan dan organisasi K3;
		\item tersedia sumber daya dan anggaran;
		\item pekerja memperoleh informasi dan pelatihan;
		\item risiko pekerjaan diidentifikasi dan dikendalikan;
		\item alat dan fasilitas kerja memenuhi persyaratan;
		\item prosedur kerja aman tersedia \& APD dikendalikan;
		\item kondisi darurat dipersiapkan;
		\item kecelakaan dan insiden diselidiki;
		\item kinerja K3 dipantau dan diperbaiki.
	\end{itemize}

	\begin{block}{Prinsip Manajerial}
		Tanggung jawab K3 berada pada level manajemen, bukan dibebankan hanya kepada petugas K3. PP No. 50 Tahun 2012 menempatkan penetapan kebijakan K3 sebagai tanggung jawab pengusaha.
	\end{block}
\end{frame}

% ---------------------------------------------------
% FRAME 11: Bukti Praktis Penerapan K3
% ---------------------------------------------------
\begin{frame}[mytitle=standard,c,mycolor=white]
	\frametitle{Bukti Praktis Penerapan K3 dalam Proyek}

	K3 harus hadir dalam dokumen dan proses pengadaan, bukan hanya di lapangan.

	\begin{block}{Dalam dokumen proyek, terdapat contoh bahwa:}
		\begin{itemize}
			\item identifikasi bahaya dan penetapan tingkat risiko telah dilakukan;
			\item persyaratan petugas/ahli K3 telah ditetapkan;
			\item setiap proses kerja dilengkapi prosedur kerja;
			\item perlindungan pekerja dicantumkan dalam proses kerja.
		\end{itemize}
	\end{block}

	\begin{exampleblock}{Makna Akademiknya}
		K3 harus terintegrasi sejak: \\
		\textbf{DED $\rightarrow$ spesifikasi $\rightarrow$ metode kerja $\rightarrow$ pemilihan penyedia $\rightarrow$ kontrak $\rightarrow$ pelaksanaan $\rightarrow$ pengawasan $\rightarrow$ evaluasi}\\
		Bukan sekadar ditambahkan setelah pekerjaan dimulai.
	\end{exampleblock}
\end{frame}

% ---------------------------------------------------
% FRAME 12: Apa Itu SMK3?
% ---------------------------------------------------
\begin{frame}[mytitle=center,c,mycolor=digiPH_leaf!10]
	\frametitle{Apa Itu SMK3?}

	\begin{block}{Definisi}
		Sistem Manajemen Keselamatan dan Kesehatan Kerja (SMK3) adalah bagian dari sistem manajemen perusahaan yang digunakan untuk mengendalikan risiko K3 secara sistematis agar tercipta tempat kerja yang aman, efisien, dan produktif.
	\end{block}

	\begin{columns}[t]
		\begin{column}{0.48\textwidth}
			\textbf{SMK3 bukan:}
			\begin{itemize}
				\item sekadar memakai helm;
				\item sekadar membuat dokumen;
				\item sekadar inspeksi;
				\item sekadar mengejar sertifikat.
			\end{itemize}
		\end{column}

		\begin{column}{0.48\textwidth}
			\textbf{SMK3 adalah:}
			\begin{itemize}
				\item sistem manajemen yang mengintegrasikan K3 ke dalam proses pengambilan keputusan dan operasi perusahaan.
			\end{itemize}
		\end{column}
	\end{columns}

	\vspace{0.5cm}
	\textit{PP No. 50 Tahun 2012 menjadi dasar utama penerapan SMK3 di Indonesia.}
\end{frame}

% ---------------------------------------------------
% FRAME 13: Tujuan dan Manfaat SMK3
% ---------------------------------------------------
\begin{frame}[mytitle=standard,t,mycolor=white]
	\frametitle{Tujuan dan Manfaat SMK3}

	\begin{block}{Tujuan}
		\begin{itemize}
			\item Meningkatkan efektivitas perlindungan K3.
			\item Mencegah dan mengurangi kecelakaan kerja serta PAK.
			\item Menciptakan tempat kerja yang aman dan produktif.
			\item Memastikan kewajiban hukum dilaksanakan.
			\item Mendorong perbaikan kinerja K3 secara berkelanjutan.
		\end{itemize}
	\end{block}

	\begin{block}{Manfaat bagi Perusahaan}
		\textbf{Manusia} $\rightarrow$ perlindungan pekerja meningkat.\\
		\textbf{Operasional} $\rightarrow$ gangguan pekerjaan berkurang.\\
		\textbf{Finansial} $\rightarrow$ biaya kecelakaan dan kerusakan dapat ditekan.\\
		\textbf{Hukum} $\rightarrow$ kepatuhan meningkat.\\
		\textbf{Reputasi} $\rightarrow$ kepercayaan pengguna jasa dan stakeholders meningkat.
	\end{block}
\end{frame}

% ---------------------------------------------------
% FRAME 14: SMK3 vs SMKK
% ---------------------------------------------------
\begin{frame}[mytitle=center,t,mycolor=digiPH_gray!20]
	\frametitle{SMK3 vs SMKK: Jangan Disamakan}

	\begin{columns}[t]
		\begin{column}{0.48\textwidth}
			\begin{block}{SMK3}
				\textbf{Cakupan:} Umum ketenagakerjaan/perusahaan\\
				\textbf{Dasar:} PP No. 50 Tahun 2012\\
				\textbf{Fokus:} Sistem manajemen K3, kebijakan, perencanaan, pelaksanaan, pemantauan, evaluasi, dan peningkatan.
			\end{block}
		\end{column}

		\begin{column}{0.48\textwidth}
			\begin{block}{SMKK}
				\textbf{Cakupan:} Khusus penyelenggaraan jasa konstruksi\\
				\textbf{Dasar:} Permen PUPR No. 10 Tahun 2021\\
				\textbf{Fokus:} Keselamatan dalam penyelenggaraan konstruksi, pengendalian risiko, integrasi keselamatan.
			\end{block}
		\end{column}
	\end{columns}

	\vspace{0.2cm}
	\begin{alertblock}{Hubungan}
		SMKK bukan pengganti prinsip K3, melainkan pendekatan sektoral untuk keselamatan konstruksi.
	\end{alertblock}
\end{frame}

% ---------------------------------------------------
% FRAME 15: Siklus Penerapan SMK3
% ---------------------------------------------------
\begin{frame}[mytitle=standard,c,mycolor=digiPH_uniblue!15]
	\frametitle{Siklus Penerapan SMK3}

	Berdasarkan Pasal 6 PP No. 50 Tahun 2012, SMK3 meliputi lima tahap:

	\begin{enumerate}
		\item \textbf{Penetapan Kebijakan K3}
		\item \textbf{Perencanaan K3}
		\item \textbf{Pelaksanaan Rencana K3}
		\item \textbf{Pemantauan dan Evaluasi Kinerja K3}
		\item \textbf{Peninjauan dan Peningkatan Kinerja SMK3}
	\end{enumerate}

	\begin{center}
		$\downarrow$\\
		\textit{Kembali ke kebijakan dan perencanaan}
	\end{center}

	\begin{exampleblock}{Artinya}
		SMK3 adalah siklus perbaikan berkelanjutan, bukan kegiatan satu kali.
	\end{exampleblock}
\end{frame}

% ---------------------------------------------------
% FRAME 16: Lima Tahap SMK3
% ---------------------------------------------------
\begin{frame}[mytitle=standard,t,mycolor=white]
	\frametitle{Lima Tahap SMK3: Dari Kebijakan hingga Perbaikan}

	\begin{table}
		\centering\small
		\begin{tabularx}{\textwidth}{|l|X|X|}
			\hline
			\textbf{Tahap} & \textbf{Pertanyaan kunci}           & \textbf{Contoh pada proyek jalan}        \\
			\hline
			1. Kebijakan   & Apa komitmen organisasi?            & Komitmen nol kecelakaan serius           \\
			\hline
			2. Perencanaan & Risiko apa yang harus dikendalikan? & Risiko alat berat, galian, lalu lintas   \\
			\hline
			3. Pelaksanaan & Bagaimana pengendaliannya?          & SOP, APD, rambu, barrier, briefing       \\
			\hline
			4. Pemantauan  & Apakah pengendalian efektif?        & Inspeksi, pengukuran, audit, investigasi \\
			\hline
			5. Peningkatan & Apa yang harus diperbaiki?          & Revisi metode kerja dan program K3       \\
			\hline
		\end{tabularx}
	\end{table}

	\begin{block}{Kunci: Plan $\rightarrow$ Do $\rightarrow$ Check $\rightarrow$ Improve}
		Kelima tahapan tersebut ditetapkan secara eksplisit dalam Pasal 6 PP No. 50 Tahun 2012.
	\end{block}
\end{frame}

% ---------------------------------------------------
% FRAME 17a: Studi Kasus dan Kesimpulan
% ---------------------------------------------------
\begin{frame}[mytitle=standard,t,mycolor=digiPH_creamy!15]
	\frametitle{Studi Kasus dan Kesimpulan}

	\begin{block}{Studi Kasus}
		Sebuah proyek peningkatan jalan melakukan pekerjaan: galian; pemadatan; penghamparan aspal; pengoperasian alat berat; dan pekerjaan pada lalu lintas aktif.

		\textbf{Pertanyaan analisis:}
		\begin{itemize}
			\item Apa saja bahaya utama? Bagaimana menentukan tingkat risikonya?
			\item Apa pengendalian yang harus diterapkan?
			\item Regulasi apa yang menjadi dasar?
			\item Bagaimana memastikan pengendalian tersebut efektif?
		\end{itemize}
	\end{block}

	\begin{exampleblock}{Kesimpulan}
		K3 konstruksi adalah kewajiban hukum sekaligus sistem manajemen. Proyek yang aman bukan proyek tanpa risiko, melainkan proyek yang mampu mengenali, mengendalikan, memantau, dan memperbaiki risikonya secara sistematis.
	\end{exampleblock}
\end{frame}

% ---------------------------------------------------
% FRAME 17b: Referensi dan Catatan
% ---------------------------------------------------
\begin{frame}[mytitle=standard,t,mycolor=white]
	\frametitle{Referensi dan Catatan}

	\begin{block}{Referensi Utama}
		\begin{itemize}
			\item UU No. 1 Tahun 1970 tentang Keselamatan Kerja.
			\item UU No. 2 Tahun 2017 tentang Jasa Konstruksi (diubah dgn UU 6/2023).
			\item PP No. 50 Tahun 2012 tentang Penerapan SMK3.
			\item UU No. 38/2004 tentang Jalan jo. UU No. 2/2022.
			\item PP No. 34 Tahun 2006 tentang Jalan.
			\item Permen PUPR No. 10 Tahun 2021 tentang Pedoman SMKK.
		\end{itemize}
	\end{block}

	\begin{alertblock}{Catatan Penting untuk Dosen}
		Untuk level sarjana akhir/pascasarjana, Slide 14–16 sebaiknya menjadi titik analisis utama: mahasiswa tidak hanya diminta menghafal regulasi, tetapi membedakan domain hukum K3 umum dan keselamatan konstruksi serta menerjemahkan 5 tahap SMK3 menjadi keputusan operasional proyek.
	\end{alertblock}
\end{frame}


\bibliographystyle{apalike}
\bibliography{../biblio.bib} % required for citecompletion, not work in mainfile
\end{document}

